\begin{frame}{거듭제곱법이 왜 수렴하는가 --- 바나흐 고정점 정리}
  \begin{itemize}
    \item 이 거듭제곱법의 수렴 논리는, \kb{ML2 Week 4}에서
      벨만 최적방정식의 해가 유일하게 존재함을 보일 때 쓸
      \kb{바나흐 고정점 정리}(Banach fixed point theorem)와
      \textbf{정확히 같은 논리}.
    \item ``\textbf{어떤 변환을 값이 안 바뀔 때까지 반복해서 적용하면
      결국 고정점}(fixed point)\textbf{에 도달한다}''는 같은 수학:
    \item 여기서 = ``\textbf{그래프 위 확률분포}'',
      강화학습에서는 = ``\textbf{상태의 가치함수}''.
  \end{itemize}
\end{frame}
